\documentclass[11pt,a4paper]{article}
\usepackage{fontspec,amsmath,amssymb,amsthm,microtype}
\setmainfont{lmroman10-regular.otf}[BoldFont=lmroman10-bold.otf,ItalicFont=lmroman10-italic.otf,BoldItalicFont=lmroman10-bolditalic.otf]
\usepackage[margin=25mm]{geometry}
\usepackage[unicode,hidelinks]{hyperref}
\hypersetup{pdfauthor={LangevinLean Project}}
\setlength{\emergencystretch}{3em}
\setlength{\parskip}{0.4em}
\setlength{\parindent}{0pt}
\allowdisplaybreaks[2]
\raggedbottom
\title{\texorpdfstring{Expected potential and gradient-\allowbreak{}step energies on L2}{Expected potential and gradient-step energies on L2}}
\author{LangevinLean Project}
\date{}
\begin{document}\maketitle
Draft mathematical exposition; these notes may contain mistakes.


The deterministic gradient-step and descent inequalities extend to the actual Hilbert space of square-integrable random variables. Potential integrability is proved from quadratic growth, and the lifted drift is identified almost everywhere with the pointwise gradient. These estimates retain the full gradient energy, including its sign-sensitive correction in the sampler one-step bound.


\section{\texorpdfstring{Integrate the strong gradient-\allowbreak{}step inequality}{Integrate the strong gradient-step inequality}}


Apply the integrated strong supporting-plane inequality to Y and X. Reverse the displacement Y-X to X-Y and multiply by nonnegative h. Expand the squared L2 distance after subtracting h times the lifted gradient. The mixed inner product is bounded by the potential difference and strong convexity, giving

\[\|X-h\nabla f(X)-Y\|_{L^2}^2\le(1-mh)\|X-Y\|_{L^2}^2+2h(\mathbb Ef(Y)-\mathbb Ef(X))+h^2\|\nabla f(X)\|_{L^2}^2.\]

The norms are genuine L2 norms and every expectation is finite.



Earlier results used here: \cite[Section~3.13.2]{LangevinLean2026ULASC}.

\section{\texorpdfstring{Integrate deterministic descent without dropping the score term}{Integrate deterministic descent without dropping the score term}}


The pointwise deterministic descent theorem holds outside a null set where the lifted state arithmetic and lifted gradient agree with their ordinary representatives. Both potential terms are integrable by smooth quadratic growth, and the squared score is integrable by the Lipschitz drift construction. Integrate the pointwise estimate, use linearity, and identify the integral of the squared lifted gradient with its squared L2 norm. The result is

\[2(\mathbb Ef(X-h\nabla f(X))-\mathbb Ef(X))\le-h(2-hM)\|\nabla f(X)\|_{L^2}^2.\]

Earlier results used here: \cite[Section~3.13.2]{LangevinLean2026ULASC}; \cite[Section~3.13.3]{LangevinLean2026ULASC}; \cite[Section~3.11.1]{LangevinLean2026IPLA}; \cite[Section~3.15.2]{LangevinLean2026Proximal}.

\begin{thebibliography}{ULA-LSI}
\bibitem[ULA-SC]{LangevinLean2026ULASC}
LangevinLean Project. \emph{Unadjusted Langevin under Strong Convexity}. LL-01, 2026. Draft report, content revision 6f6f50decdb6. \url{https://odakyildiz.github.io/LangevinLean/docs/pdf/UnadjustedLangevin/FullProof.pdf}.
\bibitem[IPLA-SC]{LangevinLean2026IPLA}
LangevinLean Project. \emph{Interacting Particle Langevin under Strong Convexity}. LL-03, 2026. Draft report, content revision c8b5d67ceb43. \url{https://odakyildiz.github.io/LangevinLean/docs/pdf/IPLA/FullProof.pdf}.
\bibitem[Proximal]{LangevinLean2026Proximal}
LangevinLean Project. \emph{Proximal particle Langevin}. LL-05, 2026. Draft report, content revision a6e5a2d101b5. \url{https://odakyildiz.github.io/LangevinLean/docs/pdf/ProximalLangevin/FullProof.pdf}.
\end{thebibliography}

\end{document}
